\chapter*{LISTE DES ABRÉVIATIONS}
\addcontentsline{toc}{chapter}{LISTE DES ABRÉVIATIONS}

\begin{longtable}{p{3.5cm}p{11cm}}
    \textbf{AP} & Agriculture de précision \\
    \textbf{FAO} & Organisation des Nations unies pour l'alimentation et l'agriculture \\
    \textbf{USDA} & Département de l'Agriculture des États-Unis \\
    \textbf{MO} & Matière organique \\
    \textbf{GPS} & \textit{Global Positioning System} \\
    \textbf{GSM} & Système mondial de communications mobiles \\
    \textbf{IoT} & Internet des objets \\
    \textbf{API} & \textit{Application Programming Interface} \\
    \textbf{CA} & Coque d'arachide \\
    \textbf{EGB} & Engrais granulé bio \\
    \textbf{PCA} & Poudre de coque d'arachide \\
    \textbf{CEC} & Capacité d'échange cationique \\
    \textbf{HT-PGPR} & \textit{Halotolerant Plant Growth-Promoting Rhizobacteria} \\
    \textbf{EC} & ElectroConductivité ou Conductivité électrique \\
    \textbf{N} & Azote \\
    \textbf{P} & Phosphore \\
    \textbf{K} & Potassium \\
    \textbf{GOM T} & Gombo témoin \\
    \textbf{GOM TS} & Gombo test \\
    \textbf{LAI T} & Laitue témoin \\
    \textbf{LAI TS} & Laitue test \\
    \textbf{NAV T} & Navet témoin \\
    \textbf{NAV TS} & Navet test \\
    \textbf{Plot T} & Plot témoin \\
    \textbf{Plot TS} & Plot test \\
    \textbf{RDU} & Réserve difficilement utilisable \\
    \textbf{RFU} & Réserve facilement utilisable \\
    \textbf{RU} & Réserve utile \\
    \textbf{SAD} & Système d'aide à la décision \\
    \textbf{MERISE} & Méthode d'étude et de réalisation informatique pour les systèmes d'entreprise \\
    \textbf{MCC} & Modèle conceptuel de communication \\
    \textbf{MCT} & Modèle conceptuel de traitement \\
    \textbf{MOT} & Modèle organisationnel de traitement \\
    \textbf{MCD} & Modèle conceptuel des données \\
    \textbf{MLD} & Modèle logique des données \\
    \textbf{PWA} & \textit{Progressive Web App} \\
\end{longtable}

\chapter*{LISTE DES SYMBOLES}
\addcontentsline{toc}{chapter}{LISTE DES SYMBOLES}

\begin{longtable}{p{3.5cm}p{11cm}}
    $\Delta$ & Pente de la courbe de pression de vapeur \\
    $\Delta S$ & Variation de la réserve en eau \\
    $\delta$ & Déclinaison solaire \\
    $\gamma$ & Constante psychrométrique \\
    $\omega_s$ & Angle horaire du coucher du soleil \\
    $\theta_{\mathrm{CC}}$ ou $H_{cc}$ & Humidité à la capacité au champ \\
    $\theta_{\mathrm{PFP}}$ ou $H_{pf}$ & Humidité au point de flétrissement permanent \\
    $\theta_{\text{actuel}}$ & Teneur en eau volumique actuelle du sol \\
    $D$ & Drainage profond \\
    $D_i$ & Dose d'irrigation \\
    $d_r$ & Distance relative Terre-Soleil \\
    $DPV$ & Déficit de pression de vapeur \\
    $e_a$ & Pression de vapeur réelle \\
    $e_s$ & Pression de vapeur saturante \\
    $ET_0$ & Évapotranspiration de référence \\
    $ET_c$ & Évapotranspiration culturale \\
    $ET_{c,\mathrm{aj}}$ & Évapotranspiration culturale ajustée \\
    $G$ & Flux de chaleur dans le sol \\
    $I$ & Irrigation \\
    $J$ & Numéro du jour de l'année \\
    $K_c$ & Coefficient cultural \\
    $K_{c,\mathrm{ini}}$ & Coefficient cultural initial \\
    $K_{c,\mathrm{mid}}$ & Coefficient cultural de mi-saison \\
    $K_{c,\mathrm{end}}$ & Coefficient cultural de fin de saison \\
    $K_{cb}$ & Coefficient cultural de base \\
    $K_e$ & Coefficient d'évaporation du sol \\
    $K_{Rs}$ & Coefficient d'ajustement \\
    $K_s$ & Coefficient de stress hydrique \\
    $p$ & Fraction de déplétion \\
    $P$ & Pression atmosphérique (précipitation dans le bilan hydrique) \\
    $R$ & Ruissellement \\
    $R_a$ & Rayonnement extraterrestre \\
    $R_c$ & Remontée capillaire \\
    $R_n$ & Rayonnement net \\
    $R_{nl}$ & Rayonnement net de grandes longueurs d'onde \\
    $R_{ns}$ & Rayonnement net de courtes longueurs d'onde \\
    $R_s$ & Rayonnement solaire \\
    $R_{so}$ & Rayonnement de ciel clair \\
    $RH_{\max}$ & Humidité relative maximale \\
    $RH_{\min}$ & Humidité relative minimale \\
    $S_a$ & Stock d'eau actuel \\
    $S_c$ & Seuil critique \\
    $T$ & Température moyenne journalière \\
    $T_{\max}$ & Température maximale \\
    $T_{\min}$ & Température minimale \\
    $u_2$ & Vitesse du vent à 2 m de hauteur \\
    $Z_r$ & Profondeur racinaire \\
    $z$ & Altitude du site \\
\end{longtable}